\chapter{A risk figure is not a decision}\label{ch:decision}

Central banks, statistical offices and health registries release granular data
under a legal duty of confidentiality. The technical means of protecting those
releases are well developed. Cell suppression, minimum cell counts, global
recoding, record swapping, the cell-key method, added noise, synthetic data and
formal privacy mechanisms are documented, implemented in software and taught
\citep{willenborg2001,hundepool2012,dwork2014}. An institution that wants to
protect a file has no shortage of tools. What it usually lacks is a rule that turns
a measured risk into a decision to release or to withhold, a rule that someone
else could apply to the same evidence and reach the same answer.

This chapter sets out why that rule is missing, what goes wrong without it, and
what the rest of the book puts in its place.

\section{How release decisions are made now}

In most institutions the decision is made by committee. A disclosure control
officer presents a table of risk measures for the proposed release: perhaps the
share of sample-unique records, the smallest cell count, a mean re-identification
probability. A panel weighs those figures against the public value of the data,
and a decision follows. The decision is often sound, because the people on the
panel know the data and the users. It is rarely reproducible, however, and it is
hard to audit afterwards, because the quantity that decided the matter was never
written down. That quantity is the residual risk the institution was willing to
carry for this release, against this kind of harm, on someone's authority.

Better protection does not repair the gap. An institution can adopt a mechanism
with a formal guarantee, account carefully for its privacy loss parameter
$\varepsilon$, and still be unable to say what probability of harm it accepted,
because $\varepsilon$ bounds how far an adversary's beliefs can move rather than
how often an adversary succeeds (Chapter~10 returns to this). The same is true of
a minimum cell count. A floor of ten is a setting of the mechanism. It becomes a
statement about risk only under assumptions about the adversary that the floor
itself does not record, and Chapter~\ref{ch:scenario} shows how quickly those
assumptions fail.

\section{Three objects that are routinely run together}

The difficulty comes from treating three different objects as one. The first is a
probability model for an attempted linkage: given what a declared adversary knows
and what the release shows, how likely is it that the adversary names the right
record. The second is a statistical estimate of that probability, which in all
but the simplest cases comes from simulating the attack many times and so carries
sampling error. The third is a policy decision: how much residual risk is
acceptable, for which harm, on whose authority and until when.

The first is a modelling problem, the second an inference problem and the third a
governance problem. Each needs its own discipline, and running them together
produces failures that are easy to miss because each looks reasonable from inside
one of the three. Two examples show the pattern.

The first failure concerns the estimate. Suppose a committee releases a file
whenever the estimated risk falls below its tolerance. For a release whose true
risk sits exactly at the tolerance, a sensible estimator lands above and below the
true value about equally often, so the committee releases about half the time.
The committee believes it has a rule that keeps risk below the tolerance. It has a
coin toss at the boundary, and nothing in its records shows this, because the
records hold point estimates. Chapter~4 states this as a proposition and shows
that replacing the point estimate with a one-sided upper confidence bound
controls the chance of a wrong release at a stated level.

The second failure concerns the decision. An institution that says ``we apply a
minimum cell count of ten'' has stated a protection setting and called it a
policy. It has not stated the probability of harm it accepted, so nobody can check
whether the setting achieves the policy, and nobody owns the judgement that the
residual risk is acceptable. When the data change, or the users change, there is
no record of what the setting was meant to deliver, so there is nothing to
reassess against.

\section{What a decision rule needs}

The book's answer is a written decision rule with three parts, one for each of
the three objects. The model part declares an adversary: who attacks, what it
knows, which targets it attacks and what conditioning event the risk figure
assumes. The inference part replaces the point estimate with an upper confidence
bound at a stated error probability. The governance part states a tolerance and
attaches to it the things that turn a number into a decision: an owner
who accepts the residual risk, an expiry date, the harm class the tolerance is set
for, the legal basis and the parties it protects, the access route through which
the data will be used, and, for a series of releases, the horizon over which the
risk is assessed.

A release is then admissible when the upper bound for every target in a declared
target set lies at or below the tolerance. This is a simple rule, and its
simplicity is the point. Two analysts given the same declarations and the same
simulation output reach the same decision, and an auditor can check afterwards
that the decision followed from the evidence.

The rule has a price, and the book does not hide it. The evidence needed to
certify a tolerance grows as the tolerance shrinks. If an institution simulates an
attack many times and the attack never succeeds, the number of clean attempts
needed for a 95 per cent upper bound to reach a tolerance of one in a hundred is
299. For one in a thousand it is 2,995, and for one in a hundred thousand it is
299,572. Observing no successful attack never shows that the risk is zero, and
choosing a tolerance and choosing the evidence standard turn out to be one decision. Chapter~4 derives these figures and discusses what an institution
should do when the price is too high.

\section{Several releases from one register}

Single releases are not the operative problem. Institutions publish programmes: a
quarterly series, a research file and a public extract, all drawn from one
register. An adversary who holds several of these is not obliged to use them one at
a time. When records keep the same pseudonym across releases, the adversary can
intersect the candidate sets that each release leaves, and the result can identify
a target that no single release exposes.

The consequence is stronger than intuition suggests. Chapter~\ref{ch:programmes}
proves that no function of the per-release risks can bound the risk of an adversary
who pools the releases. Each release can be made as safe as one likes, judged on
its own, while the pooled releases identify the target with certainty. The
quantity an institution certifies when it reviews releases one by one is therefore
not the quantity that governs its exposure. On a simulated register, four releases
that each keep at least three records in every key class, and that contain no
unique record at all, identify the target with probability 0.458 once pooled,
against an upper bound of 0.253 implied by the four marginal risks. A quarter of
targets are identified uniquely by the intersection although none was unique in
any single release.

This is the question a workshop participant raised in Lisbon, and it is the reason
the book gives programmes a part of their own.

\section{A running example}

Most chapters return to one example: a synthetic credit register built to share
the structure of a supervisory collection such as the euro area's AnaCredit
\citep{israel2017}. It holds 20,000 firms as borrowers, 60 lenders with unequal
shares, 18 sectors, 12 regions and four size classes, observed over six monthly
vintages. About 27,000 loans are live at any time. Each released loan carries a
persistent pseudonym, its sector, region and size class, and its exposure.

The register is synthetic because real loan-level credit data are not released for
general use, and that choice limits what the example can show. Its parameters are
chosen by design, so the protection settings it implies describe the simulation
and not an actual institution. What the example can show is how the decision rule
behaves when the adversary changes. At an illustrative tolerance of 0.12, an
assessment of one release at a time passes the file with no suppression. An
adversary who pools the six vintages and knows in which of them the target's loan
was live requires a minimum cell count of 75. One who also holds a reading of the
loan's exposure in each vintage requires 150, and at that floor the file keeps 20
per cent of loans but only 4 per cent of exposure, with no loan to a large firm.
Against that adversary a loan-level file is the wrong product, and aggregates or
access in a secure setting are the realistic options. Part~V develops the example
in full.

\section{How the book is arranged}

Part~I sets up the problem. Chapter~\ref{ch:scenario} defines disclosure risk as
the outcome of a declared attack, derives the familiar reciprocal-frequency
measure from four assumptions and shows what happens when two of them fail.
Chapter~3 shows that the measures reported under one label answer different
questions. Part~II treats the risk figure as an estimate: its uncertainty, the
choice among protection settings, and certification while an attack is being
monitored. Part~III turns to the decision itself: the elements of a governance
threshold, and the tolerance implied by a stated weight on harm against utility.
Part~IV deals with programmes of releases and with formal privacy. Part~V reports
the evidence from simulation, public microdata, the credit register and a pipeline
in which a model proposes attacks. Part~VI is for practice: falsification tests, an
operational protocol and worksheets for six release types.

A reader who must decide on a release next month can read this chapter, the
\textit{In practice} blocks at the end of each chapter, and Part~VI. A reader who
wants to check the claims will find each one stated as a proposition with its
assumptions, and proved in the last section of its chapter.

\begin{inpractice}
Before the next release decision, write down four things that are usually left
implicit: the adversary the assessment assumes, the tolerance the release must
meet, the person who owns that tolerance, and the date after which the assessment
lapses. If any of the four cannot be written down, the committee is not yet in a
position to say whether the release is safe, whatever the risk table shows.
\end{inpractice}
